% !TEX root = ../../tree_notes.tex
% =====================================================================
%  Chapter 16 review -- the 15-station revision contract.
%  Revision-only: tables, one-liners, pointers back via \thmref/\probref/
%  \secref/\coderef/\defref/\exref/\figref/\propref.  No derivations.
%  Summit: CF 1062E Company -- LCA of a range via RMQ on tin + binary
%  lifting.  Sample verified against the live CF statement (n=11,
%  parents 1 1 3 3 3 4 2 7 7 6; 5 queries; output "4 1 / 8 1 / 1 0 /
%  11 3 / 8 1").
%  Exports prob:c11-summit, lst:c11-summit, lst:c11-kit.
% =====================================================================

\chaptersection{Chapter 16 Summary}

One sentence: \emph{a fixed root turns every pair of nodes into a unique
hinge---the LCA---and three machines (the walk, the Euler-tour RMQ, the
binary-lifting table) answer ``where is the hinge?'' in
\(O(h)\)/\(O(1)\)/\(O(\log n)\); once you have the hinge, distances,
on-path tests, depth-bucket counts, and route-position queries all
become one arithmetic expression.} Chapter~13 made the subtree fold
queryable; Chapter~16 makes the \emph{pair} queryable. Every row below
is a pointer, not a re-derivation.

\begin{center}
\renewcommand{\arraystretch}{1.25}
\begin{tabular}{@{}p{0.22\linewidth} p{0.50\linewidth} p{0.20\linewidth}@{}}
\toprule
\textbf{Result} & \textbf{Statement} & \textbf{Where} \\
\midrule
Chain property & Common ancestors of \(u,v\) = ancestors of \(\lca(u,v)\); a single chain. & \propref{c11-chain} \\
Two-phase walk & Equalise depths, lockstep up; first agreement is \(\ell\); \(O(h)\). & \secref{c11-idea} \\
Euler--RMQ reduction & \(\lca(u,v) = \arg\min\) depth on the \(2n{-}1\) tour between \(\mathit{first}(u)\) and \(\mathit{first}(v)\); sparse table gives \(O(1)\). & \thmref{c11-eulerthm}, \defref{c11-spq} \\
Binary lifting table & \(\code{up}[v][j] = \code{up}[\code{up}[v][j{-}1]][j{-}1]\); \(O(\log n)\) climb and LCA. & \defref{c11-up} \\
Aggregated lifting & \(\code{agg}[v][j]\) carries an associative payload; path aggregate in \(O(\log n)\). & \thmref{c11-agg} \\
Distance identity & \(\dist(u,v) = \depth(u)+\depth(v)-2\depth(\ell)\). & \thmref{c11-distf} \\
On-path test & \(w \in \mathit{path}(u,v) \iff \dist(u,w)+\dist(w,v)=\dist(u,v)\). & \propref{c11-onpath} \\
Subtree interval & \(u \preceq v \iff \tin(u)\le\tin(v)\le\tout(u)\); counts = binary search. & \propref{c11-interval} \\
Depth buckets & Sorted \(\tin\) per depth; subtree count in \(O(\log n)\) via two searches. & \secref{c11-paths} \\
\bottomrule
\end{tabular}
\end{center}

\medskip
{\small\textbf{Algorithms from this chapter.}}
\begin{center}
\renewcommand{\arraystretch}{1.2}
\begin{tabular}{@{}p{0.27\linewidth} l l >{\raggedright\arraybackslash}p{0.26\linewidth}@{}}
\toprule
\textbf{Algorithm} & \textbf{Prep} & \textbf{Query} & \textbf{Code} \\
\midrule
Naive walk & \(O(n)\) DFS & \(O(h)\) & \coderef{c11-naive} \\
Euler tour + sparse table & \(O(n\log n)\) & \(O(1)\) & \coderef{c11-euler}, \coderef{c11-sparse} \\
Binary lifting (build/climb/LCA) & \(O(n\log n)\) & \(O(\log n)\) & \coderef{c11-build}, \coderef{c11-lcalift} \\
CSES 1135 (Distance Queries) & \(O(n\log n)\) & \(O(\log n)\) & \coderef{c11-1135-soln} \\
CSES 1687 (Company Queries I) & \(O(n\log n)\) & \(O(\log k)\) & \coderef{c11-1687-soln} \\
CF 1304E (route + parity) & \(O(n\log n)\) & \(O(\log n)\) & \coderef{c11-1304e-soln} \\
CF 208E (depth buckets) & \(O(n\log n)\) & \(O(\log n)\) & \coderef{c11-208e-soln} \\
Path aggregate by strides & \(O(n\log n)\) & \(O(\log n)\) & \coderef{c11-pathagg} \\
\bottomrule
\end{tabular}
\end{center}

\chaptersection{Master Decision Flowchart}

Post-study ladder: walk top to bottom, stop at the first branch that
matches. It decides \emph{which of this chapter's three LCA machines}
the problem is asking for---and hands off what this chapter does
\emph{not} own: dynamic path updates (Ch.~19/20/22 HLD).

\begin{lstlisting}[language={}, basicstyle=\ttfamily\small, numbers=none, frame=single]
  You need the LCA, or a measurement that uses it.  Which machine?
  |
  |-- h small, or q tiny (one-off / LC 236 / LC 1483)?
  |     +-- naive walk (16.1, lst:c11-naive); O(h), no build.
  |
  |-- tree STATIC, q large, ONLY need LCA (CSES 1135)?
  |     +-- Euler tour + sparse table (16.2, lst:c11-euler +
  |        lst:c11-sparse); O(1) query, idempotent combine.
  |
  |-- need k-th ancestor, distance, on-path, OR log is fast enough?
  |     +-- binary lifting (16.3, lst:c11-lcalift); O(log) per query.
  |        (Contest default: most flexible, simplest to extend.)
  |
  |-- need an AGGREGATE along the path (max/min/sum/xor)?
  |     +-- aggregated lifting (16.4, thm:c11-agg, lst:c11-pathagg).
  |        Associative op? Use two accumulators on the two arms.
  |        Idempotent? Overlapping blocks also work (sparse trick).
  |
  |-- query NAMES a SET of vertices (LCA of k nodes)?
  |     +-- RMQ-on-tin: sort by tin, LCA(set) = LCA(min-tin, max-tin).
  |        Segment tree over tin for range argmin/argmax -> Summit.
  |
  |-- values CHANGE between queries (dynamic)?
  |     +-- NOT this chapter -- HLD (Ch.~22) + segment tree (Ch.~19/20).
  |        Static tables rot; no pre-merged jump survives an update.
  |
  +-- distance DISTRIBUTION over all pairs?
        +-- NOT this chapter -- centroid (Ch.~24). Per-query LCA
            is too slow when you need O(n^2) pairs.
\end{lstlisting}

\chaptersection{Common Mistakes Checklist}

\begin{itemize}[nosep, leftmargin=1.4em]
  \item \textbf{Forgetting \code{return par} at the end of
        lifting-LCA.} The descending loop stops with \(u,v\) one level
        below \(\ell\); returning \(u\) gives a child---silent off-by-one
        on ancestor-pair tests (\coderef{c11-lcalift}).
  \item \textbf{Dropping the \code{if(u==v)} shortcut after
        equalising.} With \(u=v\) the inequality-loop never fires;
        without the shortcut the \code{par} return is wrong by one edge
        (\coderef{c11-lcalift}).
  \item \textbf{Swapping RMQ endpoints.} The span is
        \([\min(\mathit{first}(u),\mathit{first}(v)),\;
         \max(\ldots)]\); feeding \((17,14)\) instead of \((14,17)\)
        gives garbage (\thmref{c11-eulerthm}).
  \item \textbf{Overlapping blocks for non-idempotent ranges.} The
        sparse-table two-block trick requires \(x \oplus x = x\): fine
        for min/max/gcd, \emph{corrupt} for sum/xor/count
        (\defref{c11-spq}).
  \item \textbf{\(\mathit{LOG}\) one short.} Need \(2^{\mathit{LOG}} >
        n\); with \(n \le 2\cdot10^{5}\), \(\mathit{LOG}=18\) fails---19
        is the minimum. Bitloops and the table second index must both
        reach \(\mathit{LOG}\) (\defref{c11-up}).
  \item \textbf{Sentinel discipline.} Node~0 must exist in the arrays
        (\(\depth[0]=0\), \(\code{up}[0][j]=0\)); otherwise overshoots
        read a real vertex (\coderef{c11-build}).
  \item \textbf{Forests treated as trees.} CF~208E's father array has
        multiple roots: forget the super-root~0 or the \(p > \depth(v)\)
        guard and DFS misses components (\coderef{c11-208e-soln}).
  \item \textbf{Tin/tout convention cross-wired.} ``\(\tout(v) =\) last
        stamp inside \(v\)'s subtree'' and ``\(\tout(v) =\) exit-time
        tick'' are both valid---neither works mixed
        (\propref{c11-interval}).
  \item \textbf{Stack depth on a chain.} Every build here is a DFS on
        up to \(2\cdot10^{5}\) nodes; the iterative two-phase rewrite
        from Ch.~12's \coderef{c09-iter} is production-necessary.
\end{itemize}

\chaptersection{WA Clinic}

Verdicts are symptoms; the causes are the short list this chapter
drilled. Diagnose top-down---every leaf names the training problem.

\begin{figure}[H]
\centering
\begin{tikzpicture}[diagram,
    v/.style={boxnode, fill=primaryblue!12, draw=primaryblue, text width=22mm, align=center},
    b/.style={boxnode, fill=difforange!12, draw=difforange, text width=12mm},
    l/.style={boxnode, fill=successgreen!10, draw=successgreen, text width=56mm,
              font=\scriptsize, align=left}]
  \node[v] (verdict) {\textbf{Verdict on the judge}};
  \node[b, below left=8mm and 4mm of verdict] (wa) {\textbf{WA}};
  \node[b, below right=8mm and 4mm of verdict] (tle) {\textbf{TLE}};
  \node[l, below right=1mm and 2mm of wa.south] (wa1)
    {returning \(u\) instead of \code{par[u]} in lifting-LCA; dropping
     the \(u{=}v\) shortcut; off-by-one in \(\tout\) convention
     (\probref{c11-cousins}, \coderef{c11-lcalift})};
  \node[l, below left=1mm and 2mm of wa.south] (wa2)
    {RMQ span endpoints swapped; using overlapping blocks for sum/xor
     (\thmref{c11-eulerthm}, \defref{c11-spq})};
  \node[l, below=5mm of tle] (tle1)
    {naive walk on a chain \(h{=}n\); scanning subtree per query
     instead of bucket + binary search (\coderef{c11-naive},
     \probref{c11-cousins})};
  \draw[flow] (verdict) -- (wa);
  \draw[flow] (verdict) -- (tle);
  \draw[flow] (wa) -- (wa1);
  \draw[flow] (wa) -- (wa2);
  \draw[flow] (tle) -- (tle1);
\end{tikzpicture}
\caption*{WA: the hinge returns a child, or the RMQ window is inverted.
TLE: the walk pays \(O(h)\) on a chain, or the per-query scan ignores
the bucket structure. RE: \(\mathit{LOG}\) too small causes array
overflow on deep trees.}
\end{figure}

\chaptersection{Quick Reference Card}

\begin{quickref}[Chapter 16 at a glance]
\textbf{Definitions.} \(\ancestors(v)\): route to root, self included;
\(\lca(u,v)\): deepest common element; \(\lca(u,v)=u \iff u \preceq v\);
\(\code{par}(r) = 0\) sentinel.

\medskip
\textbf{Walk.} Lift deeper to equal depth, lockstep; \(O(h)\).
Check \(h\) before building tables.

\medskip
\textbf{Euler-RMQ.} \(\mathit{tour}\): enter + each return, \(2n{-}1\)
slots; \(\mathit{first}(v)\): first slot. \(\lca = \arg\min\) depth on
\([\mathit{first}(u)..\mathit{first}(v)]\). Sparse table over tour
depth: query = two overlapping \(2^j\)-windows;
\(j = \lfloor\log_2 \mathit{len}\rfloor\).
\(O(n\log n)\) prep / \(O(1)\) query.

\medskip
\textbf{Lifting.} \(\code{up}[v][j] = \code{up}[\code{up}[v][j{-}1]][j{-}1]\);
climb \(k\): jump bits low\(\to\)high. LCA: equalise, shortcut,
descend high\(\to\)low on \emph{inequality}, return \code{par}.
\(O(n\log n)\) / \(O(\log)\). Ancestry: \(\tin(a) \le \tin(v) \le
\tout(a)\).

\medskip
\textbf{Paths.} \(\dist(u,v) = \depth u + \depth v - 2\depth\ell\);
\(w \in \mathit{path}(u,v) \iff \dist(u,w)+\dist(w,v)=\dist(u,v)\);
\(k\)-th route node: climb from nearer end. Aggregates: second table
\(\code{agg}[v][j]\), merge into accumulator during climbs.

\medskip
\textbf{Translation.} ``boss \(k\) above'' \(\to\) climb.
``distance / via / route'' \(\to\) hinge identity.
``same-generation relatives under ancestor'' \(\to\) lift + depth-bucket
\(\cap\) interval (CF~208E). ``LCA of a SET of \(k\) nodes'' \(\to\)
sort by tin, LCA(min-tin, max-tin) (Summit).
\end{quickref}

\begin{lstlisting}[caption={Chapter 16 kit: the lifting machine, complete and production-shaped}, label={lst:c11-kit}]
struct LCA {   // binary lifting, 0-sentinel, iterative DFS
    int n, LOG; vector<vector<int>> adj;
    vector<array<int,20>> up;
    vector<int> dep, tin, tout; int tmr=0;
    LCA(int n):n(n),LOG(1),adj(n+1),up(n+1),dep(n+1),tin(n+1),tout(n+1){
        while((1<<LOG)<=n) LOG++;}
    void addEdge(int a,int b){adj[a].push_back(b);adj[b].push_back(a);}
    void build(int r=1){
        vector<int> par(n+1),it(n+1,0),stk{r};
        up[r].fill(0);dep[r]=0;par[r]=0;tin[r]=++tmr;
        while(!stk.empty()){int v=stk.back();
            if(it[v]<(int)adj[v].size()){int u=adj[v][it[v]++];
                if(u==par[v])continue;
                par[u]=v;dep[u]=dep[v]+1;tin[u]=++tmr;
                up[u][0]=v;for(int j=1;j<LOG;j++)up[u][j]=up[up[u][j-1]][j-1];
                stk.push_back(u);
            }else{tout[v]=tmr;stk.pop_back();}}}
    int climb(int v,int k)const{
        for(int j=0;j<LOG&&v;j++)if(k>>j&1)v=up[v][j];return v;}
    bool isAnc(int a,int v)const{
        return tin[a]<=tin[v]&&tin[v]<=tout[a];}
    int lca(int u,int v){
        if(dep[u]<dep[v])swap(u,v);
        u=climb(u,dep[u]-dep[v]);if(u==v)return u;
        for(int j=LOG-1;j>=0;j--)if(up[u][j]!=up[v][j]){u=up[u][j];v=up[v][j];}
        return up[u][0];}
    int dist(int u,int v){int w=lca(u,v);return dep[u]+dep[v]-2*dep[w];}
    bool onPath(int w,int u,int v){
        return dist(u,w)+dist(w,v)==dist(u,v);}
};
// recognition: "k above"->climb; "ancestor?/inside subtree?"->isAnc;
// "distance/via/route"->lca()+dist(); updates->NOT this kit: Ch.22.
\end{lstlisting}

\chaptersection{Write-From-Memory Drills}

Five signature skeletons, each lifted from a listing already in the
chapter. Cover the marked line, reproduce from memory, then check the
answer one line below.

\begin{lstlisting}[caption={Drill 1 -- the naive walk (from lst:c11-naive)}]
int lca_walk(int u, int v) {
    while (depth[u] > depth[v]) u = par[u];   // equalise u
    while (depth[v] > depth[u]) v = par[v];   // equalise v
    while (u != v) { u = par[u]; v = par[v]; } // lockstep
    return u;                                   // first agreement
}
\end{lstlisting}
\noindent{\scriptsize Key: no preprocessing; \(O(h)\) worst case.
Correct for small~\(h\) or tiny~\(q\); otherwise TLE on a chain.
This is the template that 16.2 and 16.3 accelerate.}

\begin{lstlisting}[caption={Drill 2 -- sparse-table RMQ query (from lst:c11-sparse)}]
int rmq_query(int l, int r) {               // min index in [l, r]
    int j = 31 - __builtin_clz(r - l + 1); // floor(log2(len))
    int a = st[j][l], b = st[j][r-(1<<j)+1];
    return (dep_tour[a] <= dep_tour[b]) ? a : b;
}
\end{lstlisting}
\noindent{\scriptsize Key: two overlapping windows of size~\(2^j\).
The overlap is harmless ONLY because min is idempotent. For sums/xors,
this trick is WRONG---use a segment tree or a disjoint sparse table
instead (\defref{c11-spq}).}

\begin{lstlisting}[caption={Drill 3 -- lifting LCA (from lst:c11-lcalift)}]
int lca_lift(int u, int v) {
    if (dep[u] < dep[v]) swap(u, v);
    u = climb(u, dep[u] - dep[v]);
    if (u == v) return u;                     // the shortcut
    for (int j = LOG-1; j >= 0; j--)
        if (up[u][j] != up[v][j]) { u = up[u][j]; v = up[v][j]; }
    return up[u][0];                          // the par
}
\end{lstlisting}
\noindent{\scriptsize Key: the loop stops with \(u,v\) as \emph{children}
of~\(\ell\). Both lines (\code{if(u==v)} and \code{return up[u][0]}) are
the two most-frequently dropped lines in contest submissions
(\coderef{c11-lcalift}).}

\begin{lstlisting}[caption={Drill 4 -- path aggregate (from lst:c11-pathagg)}]
int path_query(int u, int v, int l) {    // l = lca(u,v)
    int res = val[l];                    // fold hinge value
    int cur = u, k = dep[u] - dep[l];
    for (int j = 0; j < LOG; j++)
        if (k & (1<<j)) { res = op(res, agg[cur][j]); cur = up[cur][j]; }
    cur = v; k = dep[v] - dep[l];
    for (int j = 0; j < LOG; j++)
        if (k & (1<<j)) { res = op(res, agg[cur][j]); cur = up[cur][j]; }
    return res;
}
\end{lstlisting}
\noindent{\scriptsize Key: the hinge's own value is folded FIRST (the
two arms \emph{exclude} the landing vertex). Associativity is all
\(\oplus\) needs---no commutativity, no idempotence
(\thmref{c11-agg}, \exref{c11-aggrun}).}

\begin{lstlisting}[caption={Drill 5 -- depth-bucket query (from lst:c11-208e-soln)}]
// Query (v, p): count nodes at depth dep[v] in subtree(w) minus v
int w = climb(v, p);                     // p-th ancestor
int d = dep[v];
int lo = lower_bound(bucket[d].begin(), bucket[d].end(), tin[w]) - bucket[d].begin();
int hi = upper_bound(bucket[d].begin(), bucket[d].end(), tout[w]) - bucket[d].begin();
answer = hi - lo - 1;                    // subtract v itself
\end{lstlisting}
\noindent{\scriptsize Key: \(\mathit{bucket}[d]\) is \emph{sorted by
construction} (DFS stamps tins in order). The \(-1\) excludes~\(v\).
The guard \code{if(p > dep[v]) return 0} must precede the
\code{climb} call (\probref{c11-cousins}).}

\chaptersection{Complete Problem Checklist}

Every problem this chapter offers---4 fully solved in-section,
practice entries, one Summit, ten mixed. Tick \emph{only} after an
unaided solve; a hinted attempt counts as half.

\begin{center}
\footnotesize
\renewcommand{\arraystretch}{1.25}
\begin{tabular}{@{}c p{0.19\linewidth} l p{0.20\linewidth} p{0.24\linewidth} c@{}}
\toprule
\textbf{\#} & \textbf{Problem} & \textbf{Source} & \textbf{Key technique} & \textbf{What to look for} & \textbf{Done?} \\
\midrule
\multicolumn{6}{l}{\textbf{Section 16.2 -- Euler tour + RMQ}} \\
S16.2.1 & Distance Queries & CSES 1135 & \lv{3} Euler tour, sparse table, \(\dist\) formula & \(2n{-}1\) tour; RMQ for LCA; distance is one expression & \(\square\) \\
\midrule
\multicolumn{6}{l}{\textbf{Section 16.3 -- binary lifting}} \\
S16.3.1 & Company Queries I & CSES 1687 & \lv{3} \(k\)-th ancestor via bit decomposition & \(\code{up}\)-table; sentinel; \(-1\) on overshoot & \(\square\) \\
\midrule
\multicolumn{6}{l}{\textbf{Section 16.4 -- lifting an aggregate}} \\
S16.4.1 & 1-Trees and Queries & CF 1304E & \lv{4} distances + parity + \(k\)-th node on route & Edge used \emph{exactly once}; only \(d_2, d_3\); parity check & \(\square\) \\
\midrule
\multicolumn{6}{l}{\textbf{Section 16.5 -- distances, buckets}} \\
S16.5.1 & Blood Cousins & CF 208E & \lv{4} lifting + depth buckets + tin interval & Forest (super-root); guard; range binary search & \(\square\) \\
\midrule
\multicolumn{6}{l}{\textbf{Summit (chapter end)}} \\
S16 & Company & CF 1062E & \lv{5} RMQ-on-tin + LCA + segment tree & LCA of set = LCA(extremes by tin); remove extreme & \(\square\) \\
\bottomrule
\end{tabular}
\end{center}

\chaptersection{Recommended Study Order}

Five focused days. Day~1 internalises the hinge concept and the naive
walk; day~2 builds the \(2n{-}1\) tour and sparse table; day~3
constructs the \(\code{up}\)-table and solves lifting problems; day~4
extends to aggregates and routes; day~5 is the Summit.

\begin{center}
\small
\renewcommand{\arraystretch}{1.25}
\begin{tabular}{@{}llp{0.55\linewidth}@{}}
\toprule
\textbf{Day} & \textbf{Focus} & \textbf{Tasks} \\
\midrule
1 & Read + do 16.1 & verify \exref{c11-census} by hand on~\(T\); write \coderef{c11-naive} cold; LC~236 (Ch.1 revisit); trace \(\lca(8,10)\) walk \\
2 & Read + do 16.2 & build the \(2n{-}1\) tour table on~\(T\) by hand; verify \(\mathit{first}[\cdot]\); solve CSES~1135 end-to-end \\
3 & Read + do 16.3 & fill the \(\code{up}\)-table on~\(T\) (\exref{c11-uptab}); CSES~1687; write \coderef{c11-lcalift} from memory \\
4 & Read + do 16.4+16.5 & \exref{c11-aggrun} (max on path); CF~1304E; \exref{c11-cousinsrun} (bucket queries); CF~208E \\
5 & Summit + revise & CF~1062E; walk the decision flowchart; 10 mixed problems (first 3 minimum) \\
\bottomrule
\end{tabular}
\end{center}

\chaptersection{Difficulty Progression}

Tiers are competency gates---each row names the \emph{skill} the tier
unlocks.

\begin{center}
\small
\renewcommand{\arraystretch}{1.3}
\begin{tabular}{@{}lp{0.34\linewidth}p{0.38\linewidth}@{}}
\toprule
\textbf{Tier} & \textbf{Problems} & \textbf{What passing this tier proves} \\
\midrule
\lv{2}--\lv{3} \ \(\bigstar\bigstar\bigstar\)
  & LC~236, CSES~1135, CSES~1687
  & you can build the tour table or the \(\code{up}\)-table from a DFS, and answer \(\lca\)/\(k\)-th-ancestor without hints \\
\lv{4} \ \(\bigstar\bigstar\bigstar\bigstar\)
  & CF~1304E, CF~208E
  & you combine lifting + distances + interval queries; recognise when a path question needs the hinge split \\
\lv{5} \ \(\bigstar\bigstar\bigstar\bigstar\bigstar\)
  & CF~1062E (Summit), CF~587C, CF~629E
  & you layer RMQ-on-tin + range-argmin/max + LCA on a set; handle forest input + edge-case guards \\
\bottomrule
\end{tabular}
\end{center}

\chaptersection{Reduction Gallery}

Every entry is ``one of the three machines + one twist.'' Each line
collapses a problem family onto a listing already in the chapter.

\begin{itemize}[nosep, leftmargin=1.4em]
  \item ``LCA of a SET of~\(k\) nodes'' (CF~1062E, CSES~1688)
        \(\to\) sort by~\(\tin\); the set's LCA is the LCA of the nodes
        with minimum and maximum \(\tin\)
        (\propref{c11-interval} as the proof key).
  \item ``Distance \(a\)-\(b\) with one edge MUST be used'' (CF~1304E)
        \(\to\) three candidate distances + parity check
        (\thmref{c11-distf}, \probref{c11-1304e}).
  \item ``\(k\)-th ancestor / boss \(k\) above'' (CSES~1687, LC~1483)
        \(\to\) bit-decomposition climb
        (\coderef{c11-climb}, \defref{c11-up}).
  \item ``Count same-depth nodes in ancestor's subtree'' (CF~208E)
        \(\to\) climb + bucket + binary search
        (\coderef{c11-208e-soln}, \propref{c11-interval}).
  \item ``Max/min/sum/xor along a path'' (CSES~2134 static)
        \(\to\) aggregated lifting, two arms merge at hinge
        (\thmref{c11-agg}, \coderef{c11-pathagg}).
  \item ``\(k\)-th node on route \(u \to v\)'' (CF~587C preview)
        \(\to\) distance + hinge split + \(k\)-th ancestor
        (\propref{c11-onpath}, \secref{c11-aggregates}).
\end{itemize}

\chaptersection{Mixed Problem Set}

Ten entries from CF~1800 to CF~2400. Each refuses to say ``LCA''
directly---recognise the hinge from the \emph{measurement} the problem
requests.

\begin{mixedset}
  \practice{Tree Reconstruction}{CF 178B3 (Div.~3 E)}
    {https://codeforces.com/contest/178/problem/B3}{CF 1800 L3}
    {Multiple DFS-order queries on a static tree; identify parent-child
     from \(\tin/\tout\) range containment---this chapter's
     \propref{c11-interval} as a decision procedure.}
  \practice{Tree}{CF 191C}
    {https://codeforces.com/problemset/problem/191/C}{CF 1900 L3}
    {Paths on a tree: given~\(k\) routes \((a_i,b_i)\), count how many
     times each edge is traversed. Difference at \(\ell_i = \lca(a_i,b_i)\);
     one subtree-sum DFS finishes in \(O(n{+}k)\)
     (\thmref{c11-distf}).}
  \practice{Wizards and Roads}{CF 16D}
    {https://codeforces.com/problemset/problem/16/D}{CF 1900 L4}
    {Detect bridges via DFS lowlink; the ``is this tree edge
     bridgeable?'' test uses \(\tin/\tout\) interval containment
     exactly as \propref{c11-interval} does.}
  \practice{Alice and the Tree}{CF 1702G2}
    {https://codeforces.com/problemset/problem/1702/G2}{CF 2000 L4}
    {XOR of values along a path \(u \to v\): compute via
     \(\mathit{xor}(\mathit{root},u) \oplus
     \mathit{xor}(\mathit{root},v) \oplus \mathit{val}(\ell)\)---
     \thmref{c11-agg} with \(\oplus = \mathrm{xor}\).}
  \practice{A and B and Team Training}{CF 519E}
    {https://codeforces.com/problemset/problem/519/E}{CF 2100 L4}
    {``How many pairs have the same \(k\)-th ancestor?''---depth-bucket
     counting via lifting, exactly the CF~208E pattern
     (\coderef{c11-208e-soln}) with the bucket on both sides.}
  \practice{Minimum Incompatibility}{CF 609E}
    {https://codeforces.com/problemset/problem/609/E}{CF 2100 L4}
    {Second-MST via max-on-path: add edge \((u,v)\) to the MST, find
     the heaviest edge on the \(u \to v\) path (\thmref{c11-agg} with
     \(\max\)), swap.}
  \practice{Persistent Bookcase}{CF 466E}
    {https://codeforces.com/problemset/problem/466/E}{CF 2100 L5}
    {LCA queries on a tree built incrementally; offline queries reduce
     to \coderef{c11-lcalift} after the tree is complete.}
  \practice{Duff in the Army}{CF 587C}
    {https://codeforces.com/problemset/problem/587/C}{CF 2200 L5}
    {Shortest path through army bases on a tree: LCA + distance
     (\thmref{c11-distf}) + \(k\)-th-node-on-route
     (\coderef{c11-climb}).}
  \practice{Tree Queries}{CF 629E}
    {https://codeforces.com/problemset/problem/629/E}{CF 2300 L5}
    {Ancestor-path queries: ``is~\(u\) an ancestor of~\(v\)?'' +
     ``path passes through which special nodes?''
     \(\tin/\tout\) interval + bucket counting
     (\propref{c11-interval}).}
  \practice{Blood Cousins Returns}{CF 246E}
    {https://codeforces.com/problemset/problem/246/E}{CF 2400 L5}
    {Extended cousin counting with multiple depth constraints: the
     CF~208E bucket+climb pattern combined with DSU-on-tree for
     per-node answers (\probref{c11-cousins}, Ch.~17 forward-cite).}
\end{mixedset}

\chaptersection{The Summit -- CF 1062E Company}

\begin{notebox}[Why this is the chapter top rope]
Sections 16.1--16.5 each solve one query per pair~\((u,v)\). The Summit
asks a \emph{set} question: given a range~\([l,r]\) of \(r{-}l{+}1\)
vertices, which single vertex to \emph{remove} so the LCA of the
remaining has maximum depth? The key reduction uses this chapter's
\(\tin/\tout\) machinery: \(\lca(\text{set}) =
\lca(\text{min-tin},\text{max-tin})\), so the problem becomes ``find
the two extreme \(\tin\)-vertices in a range, then test removing each
extreme.'' The segment tree for range-argmin/argmax is \chref{c13}'s tool;
the LCA inside it is \emph{this} chapter's.
\end{notebox}

\solvedproblem{c11-summit}{Company}{Codeforces 1062E}
  {https://codeforces.com/problemset/problem/1062/E}{CF 2300 L5}
  {LCA of a set via tin-extremes, segment tree for range argmin/argmax, binary lifting, remove-one optimization}

\begin{problemstatement}[CF 1062E Company]
A company of~\(n\) employees forms a rooted tree (employee~1 is the
CEO, root).~\(q\) queries: each gives a range~\([l,r]\)---employees
\(l, l{+}1, \ldots, r\) are involved. The project manager is the
deepest common ancestor (= LCA) of all involved employees. Remove
\emph{exactly one} employee from the range so that the LCA of the
remaining employees is as \emph{deep} as possible. Output which
employee to remove and the maximum depth achievable.

\textbf{Input.} \(n\)~\(q\); then \(n{-}1\) integers
\(p_2, p_3, \ldots, p_n\) (the boss of each employee); then \(q\) lines
\((l_i, r_i)\). \(2 \le n, q \le 10^5\).

\textbf{Example.}
\begin{verbatim}
Input:
11 5
1 1 3 3 3 4 2 7 7 6
4 6
4 8
1 11
9 11
8 11

Output:
4 1
8 1
1 0
11 3
8 1
\end{verbatim}
\end{problemstatement}

\begin{attemptgate}[Attempt it yourself first -- \lv{5}, 45 minutes]
\peek{How do you compute the LCA of a SET of~\(k\) nodes? Hint:
sort them by~\(\tin\); what happens if you take the LCA of the first
and last?}
\nudge{The segment tree stores (min-tin, max-tin) in each interval.
For the ``remove one'' variant, also query for the \emph{second}
extreme: second-min and second-max. Only removing an extreme vertex
can change the LCA of the set.}
\fullhint{For each range~\([l,r]\): find \(a = \arg\min\;\tin\) and
\(b = \arg\max\;\tin\) in the range. The LCA of the whole range =
\(\lca(a,b)\). To remove one vertex: try removing~\(a\) (compute
\(\lca(\text{second-min}, b)\)) or removing~\(b\) (compute
\(\lca(a, \text{second-max})\)); pick whichever gives deeper LCA.
Output the removed vertex and the new depth. Total \(O((n+q)\log n)\).}
\end{attemptgate}

\solstep{Understanding}
The set~\(\{l, l{+}1, \ldots, r\}\) (consecutive vertex indices) has
its LCA determined by the two vertices with extreme DFS-entry times.
This is the ``LCA of a set'' reduction:
\[
  \lca(v_1, \ldots, v_k)
  \;=\; \lca\!\bigl(\arg\min_i\;\tin(v_i),\;
                     \arg\max_i\;\tin(v_i)\bigr).
\]
The proof uses \propref{c11-interval}: if~\(u\) has the smallest
\(\tin\) and~\(w\) the largest, then every other~\(v_i\) has
\(\tin(u) \le \tin(v_i) \le \tin(w)\), so all~\(v_i\)'s are in
\(\lca(u,w)\)'s subtree. Removing one vertex changes the extremes only
if that vertex \emph{is} an extreme---hence try removing the min-tin
or the max-tin vertex.

\solstep{Key Observations}
\begin{observation}
The LCA of a set = LCA of the min-tin and max-tin elements. This
turns a \(k\)-way query into one pair query, answering in \(O(\log n)\)
via binary lifting.
\end{observation}
\begin{observation}
Removing a non-extreme vertex does NOT change the LCA (the extremes
still bracket the same set). So the optimal removal is always one of
the two extremes.
\end{observation}
\begin{observation}
Second-min and second-max \(\tin\) in the range give the two candidate
LCAs after removal. A segment tree storing (min, max, second-min,
second-max) in each node answers all four in \(O(\log n)\) per range
query.
\end{observation}

\solstep{Think It Through}
\begin{thinkbox}
\deadend{Compute LCA of all~\(k\) vertices by pairwise LCA iteration}
{Correct for one query, but \(O(k)\) per query; \(O(nq)\) worst case.}
\deadend{Sort by tin, take first and last---no segment tree}
{Works only if the set is pre-sorted; here the set changes per query
(range of indices, not DFS order). You need a \emph{range} data
structure.}
\aha{Segment tree over~\(\{1, \ldots, n\}\) storing min/max/second-
min/second-max of~\(\tin[\cdot]\). Per query: four values from the
tree in \(O(\log n)\); one LCA call in \(O(\log n)\). Total
\(O((n+q)\log n)\).}
\end{thinkbox}

\solstep{Implementation}
\begin{lstlisting}[caption={CF 1062E: segment tree on tin + binary lifting LCA}, label={lst:c11-summit}]
#include <bits/stdc++.h>
using namespace std;
const int MX=100002, LOG=18;
vector<int> adj[MX];
int up[MX][LOG], dep[MX], tin[MX], tout[MX], tmr;
void dfs(int v,int p){
    tin[v]=++tmr; up[v][0]=p;
    for(int j=1;j<LOG;j++) up[v][j]=up[up[v][j-1]][j-1];
    for(int c:adj[v]){dep[c]=dep[v]+1;dfs(c,v);}
    tout[v]=tmr;
}
int lca(int u,int v){
    if(dep[u]<dep[v])swap(u,v);
    int d=dep[u]-dep[v];
    for(int j=0;j<LOG;j++)if(d&(1<<j))u=up[u][j];
    if(u==v)return u;
    for(int j=LOG-1;j>=0;j--)if(up[u][j]!=up[v][j]){u=up[u][j];v=up[v][j];}
    return up[u][0];
}
struct Node{int mn1,mn2,mx1,mx2;};
int n;
vector<int> rk; // rk[tin[i]]=i
Node merge_nodes(Node a, Node b){
    int c[4]={a.mn1,b.mn1,a.mn2,b.mn2};
    sort(c,c+4,[](int x,int y){return tin[x]<tin[y];});
    int d[4]={a.mx1,b.mx1,a.mx2,b.mx2};
    sort(d,d+4,[](int x,int y){return tin[x]>tin[y];});
    return {c[0],c[1],d[0],d[1]};
}
Node tr[4*MX];
void build(int o,int l,int r){
    if(l==r){tr[o]={l,l,l,l};return;}
    int m=(l+r)/2;build(2*o,l,m);build(2*o+1,m+1,r);
    tr[o]=merge_nodes(tr[2*o],tr[2*o+1]);
}
Node query(int o,int l,int r,int ql,int qr){
    if(ql>r||qr<l) return {0,0,0,0};
    if(ql<=l&&r<=qr) return tr[o];
    int m=(l+r)/2;
    return merge_nodes(query(2*o,l,m,ql,qr),query(2*o+1,m+1,r,ql,qr));
}
int main(){
    int q; scanf("%d%d",&n,&q);
    for(int i=2;i<=n;i++){int p;scanf("%d",&p);adj[p].push_back(i);}
    tin[0]=1e9; dep[1]=0; dfs(1,0); build(1,1,n);
    while(q--){
        int l,r; scanf("%d%d",&l,&r);
        if(l==r){printf("%d %d\n",l,dep[1]);continue;}
        Node nd=query(1,1,n,l,r);
        int mn=nd.mn1,mx=nd.mx1;
        int cmn=nd.mn2,cmx=nd.mx2;
        int lca_a=lca(cmn,mx); // remove mn
        int lca_b=lca(mn,cmx); // remove mx
        if(dep[lca_a]>=dep[lca_b])
            printf("%d %d\n",mn,dep[lca_a]);
        else
            printf("%d %d\n",mx,dep[lca_b]);
    }
}
\end{lstlisting}

\solstep{Trace on the Official Sample}
\begin{tracebox}[CF 1062E: n=11, parents 1 1 3 3 3 4 2 7 7 6; 5 queries]
Tree: \(1\!\to\!\{2,3\}\), \(3\!\to\!\{4,5,6\}\), \(4\!\to\!\{7\}\),
\(2\!\to\!\{8\}\), \(7\!\to\!\{9,10\}\), \(6\!\to\!\{11\}\).
Depths: \(d=(0,1,1,2,2,2,3,2,4,4,3)\).

Q1\;(4,6): set~\(\{4,5,6\}\). All have parent~3.
\(\lca(4,5,6)=3\), depth~1. Remove any: remaining LCA still~3.
Output \texttt{4 1}. \checkmark

Q2\;(4,8): set~\(\{4,5,6,7,8\}\).
\(\lca(4,5,6,7,8)=1\) (node~8 in subtree~2, rest in subtree~3).
Remove~8: \(\lca(4,5,6,7)=3\), depth~1. Remove others: LCA=~1.
Best: remove~8, depth~1. Output \texttt{8 1}. \checkmark

Q3\;(1,11): set~\(=\{1,\ldots,11\}\). Node~1 is in the set,
so \(\lca(\text{all})=1\), depth~0. Remove any vertex leaves~1.
Output \texttt{1 0}. \checkmark

Q4\;(9,11): set~\(\{9,10,11\}\).
\(\lca(9,10)=7\), \(\lca(9,11)=\lca(10,11)=3\).
Remove~11: LCA(9,10)~=~7, depth~3. Remove~9 or~10: LCA~=~3, depth~1.
Best: remove~11, depth~3. Output \texttt{11 3}. \checkmark

Q5\;(8,11): set~\(\{8,9,10,11\}\).
\(\lca(8,9,10,11)=1\), depth~0.
Remove~8: \(\lca(9,10,11)=3\), depth~1.
Remove others: LCA stays~1, depth~0.
Best: remove~8, depth~1. Output \texttt{8 1}. \checkmark
\end{tracebox}

\solstep{Complexity}
\begin{complexitybox}
Preprocessing: \(O(n\log n)\) (up-table via one DFS) + \(O(n)\)
(segment-tree build). Per query: \(O(\log n)\) (segment tree gives
four extremes) + \(O(\log n)\) (two LCA calls). Total
\(O((n+q)\log n)\). Memory: \(O(n\log n)\) for lifting + \(O(n)\) for
segment tree. At \(n,q \le 10^5\): roughly 20\,MB. Verdict: \ac.
\end{complexitybox}

\solstep{Pitfalls and Variants}
\begin{warningbox}[The ``remove one'' analysis]
Only removing a \(\tin\)-\emph{extreme} vertex can change the set LCA.
If the range has \(r{-}l{+}1 = 2\) (just two vertices), removing either
leaves a singleton whose LCA is itself---output the correct vertex.
\end{warningbox}
\begin{warningbox}[The segment-tree identity for empty entries]
The merge function uses~\(\tin[0]\) for ``empty'' entries when a range
query spans fewer than four vertices. Set~\(\tin[0] = 10^9\) so the
dummy never wins an argmin. Forgetting this gives RE or WA on
singleton ranges.
\end{warningbox}

\begin{insightbox}[Key takeaways --- three things to remember]
\begin{enumerate}[nosep, leftmargin=1.6em]
  \item \emph{LCA of a set = LCA of the two extremes by~\(\tin\).}
        The proof is \propref{c11-interval}: the min-tin and
        max-tin bracket the entire set's DFS window, so their common
        ancestor covers everyone (\probref{c11-summit}).
  \item \emph{Removing a non-extreme preserves the LCA.} The extremes
        are the only ``bottleneck'' vertices; the second-extreme gives
        the post-removal LCA without rescanning.
  \item \emph{This chapter's machine + \chref{c13}'s segment tree =
        set-LCA queries in~\(O(\log n)\).} The handoff: this chapter
        owns the LCA + tin/tout; \chref{c13} owns the range-min structure.
\end{enumerate}
\end{insightbox}

\begin{hardeningbox}
\begin{itemize}[nosep, leftmargin=1.4em]
  \item \textbf{Generalises}: any ``aggregate over a set of vertices
        via LCA'' (LC~1483, CF~1702G2, CF~629E) uses the same
        extremes-by-tin reduction; the combiner (max, xor, count)
        rides on the aggregate table from \thmref{c11-agg}.
  \item \textbf{Chapter~16's five sections form one pipeline}: DFS
        \(\to\) tin/tout/depth \(\to\) up-table \(\to\) LCA \(\to\)
        distance/path/aggregate/set. The Summit is the first problem
        that requires \emph{all five} layers simultaneously.
  \item \textbf{Alternatives}: Euler-tour RMQ (\thmref{c11-eulerthm})
        gives \(O(1)\) LCA inside the segment tree---slightly faster
        at the cost of a second sparse-table build.  A virtual tree
        (Ch.~26) solves the same set queries but is overkill for
        single-range problems.
\end{itemize}
\end{hardeningbox}

\chaptersection{Tool Kit Inventory}

The machine assembled through Chapter~16. Chapters~1--15's rows are
carried forward; this chapter's are appended. Pointers only.

\begin{center}
\small
\renewcommand{\arraystretch}{1.25}
\begin{tabular}{@{}p{0.21\linewidth} l c p{0.25\linewidth} p{0.17\linewidth}@{}}
\toprule
\textbf{Tool} & \textbf{Since} & \textbf{Lv} & \textbf{Trigger phrase} & \textbf{Where} \\
\midrule
\multicolumn{5}{l}{\itshape Chapters 1--15 -- carried forward} \\
Rooted-view DFS & 12 & \lv{2} & ``how deep / tall / many below'' & \coderef{c09-dfs3} \\
Iterative two-phase DFS & 12 & \lv{3} & ``chain may be deep'' & \coderef{c09-iter} \\
\(\tin/\tout\) flattening & Ch.5 & \lv{2} & ``is~\(u\) an ancestor of~\(v\)?'' & \chref{c24} \\
Subtree fold (DP template) & 13 & \lv{3} & ``one answer per node, bottom-up'' & \thmref{c10-subtree-template} \\
Rerooting & 15 & \lv{4} & ``for every root'' & \chref{c20} \\
\midrule
\multicolumn{5}{l}{\itshape Chapter 16 -- new this chapter} \\
Naive walk & 16.1 & \lv{2} & ``LCA, few queries, small~\(h\)'' & \coderef{c11-naive} \\
Euler tour + sparse table & 16.2 & \lv{3} & ``static, huge~\(q\), \(O(1)\) LCA'' & \coderef{c11-euler}, \coderef{c11-sparse} \\
Binary lifting & 16.3 & \lv{3} & ``\(k\)-th ancestor / distance / default'' & \coderef{c11-build}, \coderef{c11-lcalift} \\
Aggregated lifting & 16.4 & \lv{4} & ``max/sum/xor along a path'' & \thmref{c11-agg}, \coderef{c11-pathagg} \\
Depth buckets & 16.5 & \lv{4} & ``same-depth in subtree'' & \coderef{c11-208e-soln} \\
Set-LCA via tin-extremes & Summit & \lv{5} & ``LCA of~\(k\) nodes / range'' & \coderef{c11-summit} \\
\bottomrule
\end{tabular}
\end{center}

\chaptersection{Related Topics}

\begin{itemize}[nosep, leftmargin=1.4em]
  \item \chref{c24} (Ch.~5, Euler Tour \& Subtree Flattening) ---
        backward: owns the entry-only \(\tin/\tout\) convention that
        \propref{c11-interval} restates; this chapter's
        \(2n{-}1\) tour (\defref{c11-tour}) is a different object,
        contrasted explicitly in~16.2.
  \item \chref{c09} (Ch.~12, Tree Metrics) --- backward: the
        rooted-view DFS \coderef{c09-dfs3} is the setup pass for every
        table here; \(\depth\) is the currency the distance formula
        spends.
  \item \chref{c10} (Ch.~13, Tree DP) --- backward:
        \thmref{c10-subtree-template} and \secref{c10-distance}
        are the fold and pair-at-LCA harvest that this chapter makes
        \emph{queryable}; the segment tree in the Summit is \chref{c13}'s
        tool.
  \item \chref{c20} (Ch.~15, Rerooting) --- backward: CSES~1136's
        ``sum of distances from every node'' combines Ch.~15's
        reroot-every-node pattern with \thmref{c11-distf}.
  \item \chref{c26} (Ch.~9, Cartesian Trees) --- the \emph{reverse}
        reduction (RMQ \(\to\) LCA \(\to\) Cartesian tree) lives
        there; this chapter owns the forward direction.
  \item Ch.~19/20/22 (segment trees, lazy, HLD) --- forward: the
        answer to ``what when the values change.''
  \item Ch.~24 (centroid decomposition) --- forward: distance
        \emph{distributions} over all pairs.
  \item Ch.~26 (virtual trees) --- forward: compresses a \(k\)-node
        set via sort-by-tin + LCA of adjacent pairs; this chapter's
        \thmref{c11-distf} and lifting are the inner loop.
\end{itemize}

\begin{lastword}[Last Word --- Chapter 16]
\emph{``Chapter~13 made the tree compute; Chapter~16 makes it answer
questions about pairs. The hinge is free: one DFS fixes it, and every
path question becomes arithmetic.''}

\medskip
The whole chapter is one object (the LCA) served by three machines
(walk, tour+RMQ, lifting) and measured by three formulas (distance,
on-path, subtree-interval). The walk is honest and cheap for
small~\(h\) (\coderef{c11-naive}); the tour converts \(\lca\) into a
range-min on a \(2n{-}1\) array (\thmref{c11-eulerthm}), and the
sparse table makes it \(O(1)\) (\defref{c11-spq}); lifting replaces
\emph{steps} by \emph{strides} (\defref{c11-up}), giving \(O(\log n)\)
per query plus the ability to carry aggregates (\thmref{c11-agg}) and
test ancestry (\propref{c11-interval}). The distance theorem
(\thmref{c11-distf}) is the simplest payoff---one subtraction after
one LCA---and it unlocks the route-parity check (CF~1304E) and the
depth-bucket count (CF~208E). The Summit layers them all: set-LCA
reduction + segment-tree RMQ + binary lifting + distance arithmetic.

\medskip
Carry three reflexes forward:
\begin{enumerate}[nosep, leftmargin=1.6em]
  \item \emph{The hinge is free information.} Any path question splits
        at~\(\ell\) into two ancestor-climbs; pay \(O(\log n)\) once,
        answer any pair query without re-DFS.
  \item \emph{Idempotence determines your data structure.} Sparse table
        (overlapping blocks) only for min/max/gcd; lifting
        (non-overlapping jumps) for sum/xor/count/associative
        (\thmref{c11-agg}).
  \item \emph{LCA of a set = LCA of the tin-extremes.} This
        reduction (proved in~\probref{c11-summit}) converts every
        ``group ancestor'' question into one pair query.
\end{enumerate}

\medskip
And three habits, drilled until boring: always verify~\(\mathit{LOG}\)
against the constraint bound before coding; check the \(u{=}v\)
shortcut and the \code{return par} in every lifting-LCA implementation;
and when a problem says ``for every pair in the subtree,'' reach for
buckets and \(\tin/\tout\) before brute-forcing---the
\probref{c11-cousins} pattern is one binary search away.

\medskip
\textbf{Next:} \chref{c28} (DSU-on-tree) takes the depth-bucket idea
and makes subtree queries with \emph{per-node answers} run in
\(O(n\log n)\) total. Ch.~19/20 (segment trees) and Ch.~22 (HLD) then
make the path queries \emph{dynamic}, answering the question this
chapter deliberately defers: what happens when the values change
between queries.
\end{lastword}
